\section{Learning a certified soft-state forecasting model}
\label{sec:method}

\subsection{A fixed raw-output forecasting task}

Let $y_t\in\mathbb R^d$ be an observed time series. A causal history is
$x_t=(y_{t-m+1},\ldots,y_t)$, optionally augmented by covariates available at
time $t$. For a finite set of horizons $\mathcal H$, the target is
$Y_t=(y_{t+h}:h\in\mathcal H)$. The history length, target transformation,
forecast origins, and horizons are part of the task definition. When comparing
different history lengths, we use the same origins, starting only after the
longest history is available. This prevents a change in evaluation law from
being confused with a change in model quality.

The main implementation fits direct forecasts at each requested horizon, sharing
the geometric representation and a joint coefficient budget. This is a
conditional time-series prediction model; it need not be a generative model for
entire trajectories. Direct fitting avoids relying on an unverified relation
between one-step feature error and the distribution encountered during a long
rollout. Iterated feature-transition models are included as separate baselines
and evaluated on exactly the same raw targets.

Development, validation, calibration, and test episodes have distinct roles.
Training estimates feature maps and decoder coefficients. Validation can choose
ordinary hyperparameters or identify one representative per model family.
Calibration supplies the confidence statements and may select among the entire
family frozen before it is observed. The final test set is used only to describe
forecasting behavior independently of those decisions. A test estimate is not
treated as an exact population risk.

\subsection{Geometric dictionaries and stable soft associations}

For dictionary block $b$, partition the training histories into $C_b$ reference
groups and fit one KAHM reconstruction model to each group, following
\citet{kumar2024kahm,kumar2025collaborative,kumar2026limits}. Let
$T_{bc}(x)\in[0,1]$ denote the resulting folding score. With fixed positive
$\omega$ and $\tau$, define
\begin{equation}
 \Phi_{bc}(x)=
 \frac{(1-T_{bc}(x)+\tau)^\omega}
 {\sum_{r=1}^{C_b}(1-T_{br}(x)+\tau)^\omega},
 \qquad \Phi_b(x)\in\Delta_{C_b}.
 \label{eq:soft-features}
\end{equation}
The implementation evaluates this expression by subtracting the maximum
log-affinity before exponentiation. A small positive denominator is not replaced
by a uniform vector. The latter operation changes the representation and can
destroy meaningful relative weights even when they are representable in log
arithmetic. A unit test contains a concrete case where the direct powers sum to
less than $10^{-12}$ but one normalized association is essentially one.

The original repository's KAHM training and folding routines are reused, with
their source hashes retained in the experiment manifest. In the original
experimental campaigns, history standardization uses training moments only.
The separate follow-up in Section~\ref{sec:input-scaling} compares this
preprocessing with training-only range scaling and bounded raw histories.
Geometry-fitting input subsampling and the maximum
number of local training points are fixed in the configuration. The resulting
maps are then held fixed when subsequent decoder fits and calibration are
performed.

Retain the stack $\Gamma_s(x)=(\Phi_1(x),\ldots,\Phi_s(x))$. Each block lies in
a simplex; the concatenation lies in a product of simplices. Different blocks
need not have nested cluster partitions. It is the retention of their actual
coordinate functions that preserves the old forecast class. Our experiments
use $C_b\in\{2,4,8,16\}$, so the four retained feature dimensions are
$2,6,14,30$.

\subsection{A constrained multihorizon decoder}

Collect the horizon targets as rows of a matrix $Y$. For a retained stack with
training feature matrix $F=[F_1,\ldots,F_s]$, fit
\begin{equation}
 \min_{a,W}\ \frac1N\|Y-\mathbf1a^\top-FW\|_{\rm F}^2
                   +\lambda\|W\|_{\rm F}^2
 \quad\text{subject to}\quad
                \sum_{b=1}^s\|W_b\|_{\rm F}\leq\kappa.
 \label{eq:training-objective}
\end{equation}
Here $W_b$ contains the rows associated with block $b$ and all output horizons;
the intercept is unpenalized. We use row-vector notation in
\eqref{eq:training-objective}; the equivalent column-vector decoder in the
theory section transposes these matrices. The complete prediction is the
coordinatewise projection of $a+W^\top\Gamma_s(x)$ onto the task's known
output box. The objective is convex before output clipping. The clipping step
is used at prediction time and is included in every reported loss.

Centering eliminates the intercept. Projected accelerated gradient then solves
the constrained quadratic. Projection is inexpensive: compute the Frobenius
norm of each coefficient block, project that nonnegative vector onto an
$\ell_1$ ball of radius $\kappa$, and rescale the blocks. The implementation
records a convex first-order optimality gap, not merely the change between
successive iterates. Appendix~\ref{app:optimization} gives the projection and
gap derivations. Accurate optimization is useful for interpreting the learner;
the held-out fitted-risk certificate remains valid even if the solver stops
early, provided the returned coefficient budget is feasible.

The ridge term is a numerical and statistical regularizer, whereas the block
constraint defines the sensitivity class. Exact replication leaves the
constrained class unchanged, but splitting coefficients over duplicate blocks
can change a squared-Frobenius ridge penalty. Replication-invariance claims below
concern feasible prediction classes and the obstruction. A separate zero-ridge
numerical test checks equality of duplicate and nonduplicate optimum risks
within their reported optimization gaps.

\subsection{Calibration and model selection}

Algorithm~\ref{alg:procedure} describes the implemented finite-family procedure.
The candidate collection includes history lengths, dictionary prefixes,
geometric feature types, and baselines. A representation-specific lower bound
and fitted-model upper bound use the same raw target, origin law, and decoder
sensitivity budget. An upper bound for the raw signal is never paired with a
lower bound for a different latent-coordinate target.

\begin{figure}[t]
\centering
\fbox{\begin{minipage}{0.94\linewidth}
\textbf{Algorithm 1: Certified geometric forecasting}
\begin{enumerate}[leftmargin=*,itemsep=3pt]
\item Specify horizons, causal origins, known target bounds, tolerances,
confidence level, and a candidate-cost rule.
\item On development data, fit geometric blocks, retain their prefixes, and
fit constrained decoders. Fit the comparison models on the same raw targets.
\item Freeze the entire candidate family, preprocessing, sensitivity budgets,
and any proposed refinement comparisons.
\item On new independent episodes, compute one averaged squared loss per
episode and horizon. Pair episodes disjointly and compute target-distance
obstruction scores at corresponding origins.
\item Allocate the failure probability over all requested risk statements and
improvement comparisons. Compute simultaneous lower and upper endpoints.
\item Return the least costly candidate whose upper endpoints meet all
tolerances. Also report excluded classes and inconclusive candidates. If no
candidate passes, return an unresolved decision.
\end{enumerate}
\end{minipage}}
\caption{The finite-family learner and its certification decision. A new model
fitted after inspecting calibration requires fresh evaluation data or a
separately justified adaptive guarantee.}
\label{alg:procedure}
\end{figure}

The lower endpoint can justify expanding the representation or changing the
sensitivity budget, but cannot prove that a proposed modification helps. We
therefore additionally compare old and new fitted forecasts through paired
episode losses. The procedure can certify an actual improvement even when the
class obstruction is zero. Conversely, it may exclude a restricted class even
if no tested refinement attains the desired accuracy.

In the numerical study, model cost is the number of retained feature
coordinates, with fixed tie-breaking by history length and validation error.
This is a transparent representation-size proxy, not an inference-runtime
measurement. For a deployment with unequal feature-evaluation costs, the same
selection theorem permits any cost function fixed before calibration.
